\documentclass[12pt]{article}
\usepackage[T1]{fontenc}
\usepackage[utf8]{inputenc}
\usepackage{lmodern}
\usepackage{textcomp}
\usepackage{enumitem}
\usepackage{geometry}
\geometry{margin=1in}
\begin{document}
\begin{center}
Answer Key - Chemistry - Undergraduate Level Assessment 4 \
\small (Answers are ordered exactly as in the question paper.)
\end{center}

\section*{Multiple Choice}
\begin{enumerate}[label=\arabic*.]
\item \textbf{Answer:} D
\\ \textit{Options:} A) BaH 2; B) BaOH; C) Ba; D) BaO
\\ \textit{Note:} The anhydride of Ba(OH)2 is BaO because the anhydride of a metal hydroxide is its corresponding metal oxide, formed by the removal of water (H$_{2}$O) from the hydroxide.
\item \textbf{Answer:} A
\\ \textit{Options:} A) Blackbody radiation curves; B) Emission spectra of diatomic molecules; C) Electron diffraction patterns; D) Temperature dependence of reaction rates
\\ \textit{Note:} Planck's quantum theory accurately describes blackbody radiation curves by introducing the concept of quantized energy levels, which resolved the ultraviolet catastrophe by predicting the intensity of radiation emitted at various wavelengths.
\item \textbf{Answer:} A
\\ \textit{Options:} A) B$_{12}$ icosahedra; B) B$_{8}$ cubes; C) B$_{6}$ octahedra; D) B$_{4}$ tetrahedra
\\ \textit{Note:} The solid-state structures of the principal allotropes of elemental boron are characterized by a network of B$_{12}$ icosahedra, which are unique polyhedral arrangements that contribute to boron's complex bonding and stability in various allotropes.
\item \textbf{Answer:} D
\\ \textit{Options:} A) F; B) Si; C) O; D) Ca
\\ \textit{Note:} Calcium has the lowest electron affinity among the given atoms because it is an alkaline earth metal with a full outer electron shell, making it less energetically favorable to gain an additional electron compared to elements with higher electron affinities.
\item \textbf{Answer:} D
\\ \textit{Options:} A) The orbitals are degenerate.; B) The set of orbitals has a tetrahedral geometry.; C) These orbitals are constructed from a linear combination of atomic orbitals.; D) Each hybrid orbital may hold four electrons.
\\ \textit{Note:} The statement is incorrect because each $sp^3$ hybrid orbital can hold a maximum of two electrons, not four, adhering to the Pauli exclusion principle.
\end{enumerate}

\section*{True / False}
\begin{enumerate}[label=\arabic*.]
\setcounter{enumi}{5}
\item \textbf{Answer:} True
\\ \textit{Options:} A) True; B) False
\\ \textit{Note:} In a three-component mixture, the maximum number of phases at equilibrium can be determined using the Gibbs phase rule, which states that the number of phases (P) is equal to the number of components (C) minus the degrees of freedom (F) plus two, and in the case of three components, it can allow for up to five phases.
\item \textbf{Answer:} True
\\ \textit{Options:} A) True; B) False
\\ \textit{Note:} The infrared absorption frequency of deuterium chloride (DCl) is shifted from that of hydrogen chloride (HCl) because the heavier reduced mass of DCl results in a lower vibrational frequency compared to HCl, reflecting the relationship between mass and vibrational motion in molecular spectroscopy.
\item \textbf{Answer:} True
\\ \textit{Options:} A) True; B) False
\\ \textit{Note:} The statement is true because the resonance frequency of 31 P nuclei in a magnetic field is determined by the Larmor equation, which shows that at a magnetic field strength of 20.0 T, the frequency is indeed 345.0 MHz.
\item \textbf{Answer:} False
\\ \textit{Options:} A) True; B) False
\\ \textit{Note:} The orbital angular momentum quantum number, l, for the outermost electron in ground-state aluminum, which has the electron configuration [Ne]3 s²3 p¹, is actually 1, corresponding to the p orbital, not 2.
\item \textbf{Answer:} True
\\ \textit{Options:} A) True; B) False
\\ \textit{Note:} Both paramagnetism and ferromagnetism arise from the alignment of unpaired electrons' magnetic moments in the presence of an external magnetic field, which is essential for these materials to exhibit magnetic properties.
\end{enumerate}

\section*{Long Form Response}
\begin{enumerate}[label=\arabic*.]
\setcounter{enumi}{10}
\item \textbf{Answer:} Lewis acids are defined as species that can accept an electron pair to form a covalent bond, often characterized by an incomplete octet or the presence of a central atom capable of expanding its valence shell. In the case of sulfur dichloride (SCl 2), the sulfur atom has a complete octet with two bonding pairs and two lone pairs, which limits its ability to accept additional electron pairs. Furthermore, SCl 2's molecular geometry is bent, leading to a lower ability to interact with electron donors, making it less electrophilic compared to other compounds like SF 4 or SO 2, which can engage more readily in Lewis acid behavior. Consequently, due to its complete octet and steric hindrance, SCl 2 is the least likely to exhibit Lewis acid behavior among similar compounds.
\\ \textit{Note:} correct because it accurately identifies that SCl 2 has a complete octet and a bent molecular geometry, which restricts its ability to act as a Lewis acid by limiting its capacity to accept additional electron pairs.
\item \textbf{Answer:} Dopants play a crucial role in modifying the electrical properties of semiconductors, and arsenic-doped silicon serves as an excellent example of an n-type semiconductor. In this process, arsenic, which has five valence electrons, is introduced into the silicon crystal lattice, which has four valence electrons. The extra electron from the arsenic atom becomes a free charge carrier, increasing the material's conductivity. This enhancement in electron concentration allows n-type silicon to effectively conduct electricity, making it essential for various electronic applications such as transistors and diodes. The ability to control the density and mobility of these electrons through doping is vital for the development of modern electronic devices.
\\ \textit{Note:} The answer correctly explains that arsenic-doped silicon functions as an n-type semiconductor by introducing extra electrons as charge carriers, thereby enhancing electrical conductivity, which is fundamental to semiconductor applications.
\end{enumerate}

\vfill
\noindent\textit{End of Answer Key}
\end{document}